\documentclass[11pt]{article}
\usepackage[margin=1in]{geometry}
\usepackage{setspace}
\usepackage{hyperref}
\hypersetup{colorlinks=true, allcolors=black}
\setlength{\parindent}{0pt}
\setlength{\parskip}{0.5em}
\setstretch{1.08}
\pagestyle{empty}

\begin{document}

\begin{center}
{\Large Teaching Statement}\\[0.3em]
Ke Cheng HSU (Clio)
\end{center}

I have taught at the Hong Kong University of Science and Technology since 2020. From 2020 to 2025 I was a teaching assistant for undergraduate macroeconomics, the economics of uncertainty and information, and Ph.D.\ microeconomic theory. I led tutorials, held office hours, graded, and for some courses wrote the tutorial material and examination questions. The average evaluation was 4.72 out of 5. In 2026 I designed and lectured the mathematics camp for the MSc in Economics. The students needed calculus, linear algebra, optimization, and dynamic programming before econometrics and theory, and many of them had not studied mathematics as undergraduates. The instructor evaluation was 4.53 out of 5. As Teaching Assistant Coordinator in the Business School in 2025--2026, I helped postgraduate students prepare teaching materials and watched their practice teaching.

I can teach business analytics and the empirical courses around it, in English, at undergraduate and master's level. That includes applied econometrics, data analysis for business decisions, and a course in which students estimate a model on real transaction or retail data rather than only fitting a prediction. My own papers use a Taobao apparel panel and Japanese bookstore data. In class I would start from a decision a firm actually faces --- whether to pay for return shipping, how to design a product that some consumers will search for --- and only then write the likelihood or the regression. I also use R for estimation, Python for data collection, C++ where a routine is slow, and Git so that a project can be rerun. I would teach those tools as part of a business-analytics or applied-machine-learning course, and I would require students to check a result against the data rather than accept code that a tool has generated.

The classroom rule is that a student should be able to say, in plain language, why the result occurs. A long queue is not proof that a shop is good. It is a signal, and later people can rationally join even when their own information says otherwise. I use that example before any formula, then ask students to work the step that is still unclear. One question usually belongs to several people. The same order --- question, mechanism, then formula --- is how I would write notes, problem sets, and examinations for a business-analytics course.

\end{document}
